\section{Reading the machine code}
\label{sec:sass}

Source level reasoning about a GPU kernel is a hypothesis about what the compiler
did. Disassembling it is the check. Every rung of the ladder, the whole tile
family, the stream-K and split-K kernels and the CUTLASS reference line are
captured to SASS under \texttt{experiments/sass/}, normalized so the captures
diff, and the top kernel's capture is a committed golden. This is the
microbenchmarking-and-disassembly method of Jia and colleagues~\cite{jia2018}
applied at a much smaller scale: not to reverse engineer the part, but to check my
own claims about my own kernels.

Everything in this section is a compile time fact. None of it needs a GPU, none of
it moves with the clock, and none of it is pending.

\subsection{The golden gate}

\texttt{scripts/sass\_diff.py} recaptures the top kernel from a fresh build and
compares it to the committed golden, exiting non-zero with a unified diff on any
difference. Ground rule 8 of this project says a gate that cannot fail is not a
gate, so both halves were demonstrated before it was trusted: rebuilding with no
source change passes, and adding one stray factor to one line of the epilogue
fails with the extra \texttt{FMUL} in the diff alongside the register allocation
that shifted around it. When a kernel change is intended the golden is regenerated
in the same commit as the kernel, so a reviewer sees exactly what moved in the
instruction stream.

The gate runs locally through \texttt{make sass-diff} and not yet in CI, for a
reason worth recording rather than papering over: the container build job runs a
13.0.1 toolkit image and this golden was assembled by nvcc 13.3.73. Two assemblers
give two answers and the job would fail on the header line before it reached an
instruction. Either the container moves to the toolkit the golden names, or the
golden is regenerated on the container's toolkit and this machine becomes the one
that disagrees. That is an ownership decision and it is open.

\subsection{What the instruction mix shows}

Counted from the committed captures, per kernel body, which for the mainloop rungs
is one K stage of the software pipeline. Three readings stand out.

The ladder's story is visible in the mix. The \texttt{mma\_ptx} rung reads
fragments scalar out of shared memory: 24 \texttt{LDS}, no \texttt{LDSM} at all,
for 8 \texttt{HMMA}. Adding \texttt{ldmatrix} replaces those with 42
\texttt{LDSM} and gets 56 \texttt{HMMA} out of the same shape. The top kernel adds
\texttt{cp.async}, and its 12 \texttt{LDGSTS} are the first instructions on this
ladder that move global data into shared memory without passing through a
register.

The barrier count collapses when the pipeline deepens. The WMMA and
\texttt{ldmatrix} rungs run 14 \texttt{BAR} per body; the top kernel and the whole
tile family run one. With three or more stages the buffer written at iteration $s$
was last read at iteration $s-1$, which the iteration-$s$ barrier already fences,
so one barrier per K stage is enough.

And the top kernel and the CUTLASS reference line amortize \texttt{ldmatrix}
identically, at 2.67 \texttt{HMMA} per \texttt{LDSM} on both sides. Whatever
separates those two kernels in wall clock, it is not how many \texttt{mma.sync}
each \texttt{ldmatrix} feeds. Their epilogues, by contrast, are opposites: the top
kernel has zero \texttt{STS} and stores 64 bit pairs straight from the fragment
layout, while CUTLASS runs 128 \texttt{STS} and 74 \texttt{LDS} to stage the
accumulator through shared memory into fully coalesced stores. Which trade wins on
this part is the first thing the gap analysis of Section \ref{sec:cutlass} should
settle.

\subsection{The register allocation study}

\texttt{benchmarks/register\_study.py} reassembles the top kernel across register
ceilings and blocks per SM floors and writes
\texttt{experiments/results/register\_study.csv}. Table \ref{tab:registers} is
generated from that file. Every column in it is a compile time fact except
throughput, which reads pending and stays that way until the owner runs it at
locked clocks.

\begin{table}[htbp]
\centering
\begin{tabular}{llrrrrrrll}
\toprule
axis & ceiling & blocks/SM floor & registers & spill st. & spill ld. & blocks by reg. & blocks/SM & limiter & GFLOP/s \\
\midrule
maxrregcount & none & 1 & 122 & 0 & 0 & 2 & 2 & shared memory & pending \\
maxrregcount & 64 & 1 & 64 & 544 & 544 & 4 & 2 & shared memory & pending \\
maxrregcount & 80 & 1 & 80 & 172 & 164 & 3 & 2 & shared memory & pending \\
maxrregcount & 96 & 1 & 96 & 12 & 12 & 2 & 2 & shared memory & pending \\
maxrregcount & 111 & 1 & 109 & 0 & 0 & 2 & 2 & shared memory & pending \\
maxrregcount & 128 & 1 & 122 & 0 & 0 & 2 & 2 & shared memory & pending \\
maxrregcount & 168 & 1 & 122 & 0 & 0 & 2 & 2 & shared memory & pending \\
launch bounds & none & 1 & 122 & 0 & 0 & 2 & 2 & shared memory & pending \\
launch bounds & none & 2 & 122 & 0 & 0 & 2 & 2 & shared memory & pending \\
launch bounds & none & 3 & 80 & 172 & 164 & 3 & 2 & shared memory & pending \\
launch bounds & none & 4 & 64 & 544 & 544 & 4 & 2 & shared memory & pending \\
\bottomrule
\end{tabular}

\caption{The register allocation study, from
\texttt{experiments/results/register\_study.csv}. Registers, spills, occupancy and
the limiter are compile time facts from ptxas; throughput is pending.}
\label{tab:registers}
\end{table}

There is a trap in running this study that is worth stating, because it costs an
hour to rediscover. \texttt{nvcc -maxrregcount} is ignored for a kernel carrying
\texttt{\_\_launch\_bounds\_\_}, and so is \texttt{ptxas -maxrregcount} when the
PTX carries \texttt{.maxntid}. Launch bounds win, by design. The kernel declares
\texttt{\_\_launch\_bounds\_\_(256)}, so a naive sweep returns the same register
count in every row and the table says nothing. The study therefore emits PTX once
with the build's own flags and edits one performance tuning directive in a copy of
it.

Three readings come out of the table, all of them compile time facts rather than
performance claims.

\textbf{109 registers are free.} A ceiling of 111 lands the kernel at 109 with no
spill at all, and 96 costs only 12 bytes each way. The kernel does not need 122
registers; ptxas simply has no reason to use fewer when nothing asks it to.

\textbf{Nothing is bought by asking.} The occupancy arithmetic on this part is
65536 registers, 100 KB of shared memory and 48 warps per SM. At 256 threads and
49152 bytes of shared memory per block, shared memory alone allows two blocks per
SM, and it stays the binding constraint at every register count in the table. At
64 registers the register file would allow four blocks; the shared memory budget
still allows two. Cutting registers buys spill traffic and no residency.

\textbf{The blocks per SM floor is the same lever from the other end.} Asking for
three blocks per SM produces exactly the same register count and spill volume that
the matching \texttt{maxrregcount} ceiling produces, because ptxas arrives there
the same way.

The practical consequence is that the occupancy lever on this kernel is the shared
memory budget and not the register budget, which is a useful thing to know before
spending a round on register pressure. What it does not tell me is whether higher
occupancy would be faster, because that is a measurement and it is pending.
